\section{Thirteen integers: a seeded guard stack and a calibrated carry block}
\label{sec:jones-seeded}

The polynomial specified below has exactly thirteen unrestricted integer
unknowns, exact total degree $57\,866\,184$, and a division-free circuit of
$759$ live binary operations ($520$ multiplications, $215$ additions and
$24$ subtractions). Its integer zero existence is equivalent to a finite
syntactic ZFC refutation. Thus, with the stated inherited compiler theorem,
ZFC proves that its lack of integer zeros is equivalent to
$\operatorname{Con}(\mathrm{ZFC})$. This is an ordinary mathematical proof,
not a new Lean certificate, a proof of consistency, or a global minimality
claim. Unlike the ten-integer companion, the arithmetic finalizer here uses
only the established Jones ratio interface, not Sun's 2026 bridge.

\subsection{A coherent fixed source and an explicit polynomial}
Let $p_1,\ldots,p_{22}$ be the original package-31 quadratic residuals after
substituting $x_i=n_i+1$ in all $31$ positive source coordinates. The exact
source is \texttt{upstream/parent\_quartic31.json}; the complete fixed port
mapping is $(\beta,K_0,M_C,M_F)=(\beta,K,M_C,M_{F,\mathrm{source}})$ from
\texttt{FIXED\_BINDING.json}. All four compressed little-endian payloads are
included and hash-authenticated. This uses the original transport coordinate,
not package 45's extra transport shift. The imported source theorem says
\[
 \exists n\in\mathbb N^{31}\ \bigwedge_{j=1}^{22}p_j(n)=0
 \quad\Longleftrightarrow\quad
 \text{there is a finite ZFC refutation}.
\]
The fifth residual equals $1$ at $n=0$, so any common zero is nonzero.

Use the selective weights $w_i$, row order, and targets $q_j$ in
\texttt{CODE.json}, identical to the companion's certified tuple.
Their maxima are $W=183$ and $Q_*=5190$; put $K_*=Q_*+2W=5556$ and
$E_*=K_*+W+2=5741$. For
$p_j(n)=\sum_\alpha p_{j,\alpha}n^\alpha$, set
\[
 c_{j,\alpha}=(2-|\alpha|)!\prod_i\alpha_i!\,p_{j,\alpha},\qquad
 D_c(B)=\sum_{j,\alpha}c_{j,\alpha}B^{q_j-\alpha\cdot w}.
\]
The finite isolation certificate checks all $71\,808$ possible interactions
of the $136$ used source terms with the $528$ padded quadratic monomials.
It proves
\[
 [t^{q_j}]\Bigl(1+\sum_i n_it^{w_i}\Bigr)^2D_c(t)=2p_j(n).
\]
It also proves that the displayed exponents in $D_c$ are distinct and its
leading coefficient is $2$. All $|c_{j,\alpha}|<A_*=2^{2134927}$;
the verifier derives this bound from the actual fixed-port bit lengths.

The unknowns are
\[
 (\epsilon,g,h,s,w,\phi,i,j,z_1,z_2,z_3,r,v)\in\mathbb Z^{13}.
\]
Every subsequent symbol is a polynomial register, not an additional unknown.
Fix
\[
 \kappa=2134940,\quad a=2^\kappa=8192A_*,\quad
 \eta=2^{16012040}=2^{15\kappa/2-10},
\]
and define
\begin{align*}
 b&=2\epsilon^2,& B&=ab^2,\\
 \ell&=\sum_iB^{w_i},& T_*&=\sum_jB^{q_j},\\
 N_0&=B^{W+1},& L&=2bB^{E_*},\\
 N&=\eta B^{2966},& c&=bB^W,\\
 M_c&=N_0-1-(b-1)\ell,&
 J_c&=2(1+g)^2D_c(B)+T_*,\\
 S_c&=g+N_0J_c,& T_c&=M_c+N_0(B-2)T_*,\\
 R&=L^2+(S_c+T_c)L+(L-1-S_c),& G_c&=32c.
\end{align*}
All powers have fixed nonnegative integer exponents. In particular, $N$ is
an ordinary polynomial, not a variable-exponent abbreviation.

Keep the literal Jones registers, with $\omega=w^2$:
\begin{align*}
 U&=N^2\omega,&Y&=N^2s,\\
 M&=RY,&P&=2M^2U,\\
 K&=R+1+h(P-1),&A&=M(U+1),\\
 B_1&=2R+1,&W_1&=b\omega,\\
 C&=B_1+W_1+\phi,&J&=4A-5,&
 Q&=3W_1C-2(W_1^2-1),\\
 D&=(A^2-1)C^2+1,&E&=iJDC^2,&F&=(A^2-1)E^2+1,\\
 G&=A+F(CD+1-A),&H&=B_1+jC,&I&=(G^2-1)H^2+1,\\
 a_0&=(P^2-1)K^2+1,&a_1&=DFI,\\
 u&=JF,&t&=QF+J(H-C).
\end{align*}
Set
\[
 \delta=C-KY,\qquad \Delta=g(c-g)K^2-c^2\delta^2,
 \qquad (d_1,d_2,d_3)=(3s-1,3\phi+2,3\Delta-1).
\]
The new seeded stack is
\[
 C_0=a_0,\qquad B_j^*=1+4d_jC_{j-1}z_j^2,\qquad
 C_j=C_{j-1}B_j^*\quad(1\le j\le3),\qquad C_*=C_3.
\]
Finally define
\begin{align*}
 Z&=\epsilon wi z_1z_2z_3a_0,&\Sigma&=Zu^5,&\rho&=t+ur,\\
 \Theta&=(\rho^3-u^2\rho-u^3)(\rho^2+23u^2),&
 X&=\Sigma v-\Theta,\\
 A_0&=\Sigma^2C_*,&A_1&=\Sigma^2C_*^2a_1.
\end{align*}
The requested thirteen-integer equation is the following single polynomial:
\begin{equation}\label{eq:jones-seeded}
 \boxed{\mathcal F_{13}=(X^2-A_0-A_1)^2-4A_0A_1=0.}
\end{equation}
The literal division-free DAG is \texttt{generated/signed\_seeded.json}.

\subsection{Sparse padding and unconditional size estimates}
Assume first only $b\ge2$. Since $B-1>A_*$ and the top coefficient of $D_c$
is $2$, we have $D_c(B)>B^{Q_*}>0$. Also $0<M_c<N_0$.
For every $g\ge1$, it follows that $S_c,T_c>0$ and
\[
 R=L^2+S_c(L-1)+T_cL+L-1>L^2.
\]
Moreover
\[
 N^2=2^{-20}a^{15}B^{5932}<B^{5947}<B^{11482}<L^2,
 \qquad b\le c\le N.
\]
Thus $R\ge N\ge8$ and $b\le N,R$ before assuming any carry or square
condition. Negative powers here occur only in this comparison proof, never
in the equation.

If $1\le g<4bB^W$, integrality gives $1+g\le4bB^W$ and $g<N_0$.
The injective coefficients imply
\[
 |D_c(B)|<\frac{A_*B^{Q_*+1}}{B-1},\qquad
 2(1+g)^2D_c(B)<\frac{32A_*b^2}{B-1}B^{K_*+1}.
\]
Because $B=8192A_*b^2$, the last prefactor is less than $1/128$.
Together with $T_*<B^{Q_*+1}/(B-1)$ this gives
$0<J_c<B^{K_*+1}$. Likewise $(B-2)T_*<B^{Q_*+1}\le B^{K_*+1}$.
Consequently
\[
 0<S_c,T_c<B^{E_*},\qquad S_c+T_c<2B^{E_*}<L.
\]
These estimates hold for arbitrary integer $b\ge2$, not just powers of two.

Now suppose $b=2^k$. Then $B=2^\nu$, where $\nu=\kappa+2k$.
The binary mask $M_c$ forces exactly
\[
 g\mathbin{\&}M_c=0\quad\Longleftrightarrow\quad
 g=\sum_i n_iB^{w_i},\qquad 0\le n_i<b.
\]
Here $\&$ is used only in the proof. For such a code, write
$h_e=[t^e](1+\sum n_it^{w_i})^2D_c(t)$. Then
\[
 |h_e|\le A_*(1+\textstyle\sum_i n_i)^2
 <1024A_*b^2=B/8<B/4.
\]
The unnormalized digits of $J_c$ are $2h_e+\varepsilon_e$, where
$\varepsilon_e=1$ at a target and $0$ elsewhere. Starting with incoming
carry $c_0=0$, normalized carries obey
\[
 c_{e+1}=\left\lfloor\frac{2h_e+\varepsilon_e+c_e}{B}\right\rfloor,
 \qquad c_e\in\{-1,0\}.
\]
At a target, its normalized digit avoids $B-2$ exactly when it is $0$ or $1$.
If $h_e=0$, the digit is $1+c_e\in\{0,1\}$. If $h_e>0$, it is at least $2$;
if $h_e<0$, a borrow normalizes it to a digit at least $2$.
Thus the target test is equivalent to $h_e=0$, independently of incoming
borrows. There is no need to fill every position with a dense $B/2$ offset.
It follows, by splitting at $N_0$, that
\[
 S_c\mathbin{\&}T_c=0\quad\Longleftrightarrow\quad
 g\text{ codes a common zero of all }p_j.
\]
Conversely every common zero has, for arbitrarily large $b=2^{1+2r}$, such a
code with $1\le g<bB^W=c$.

\subsection{The parity-tuned two-block threshold}
Write $L=2^\lambda$ after $b=2^k$ has been established. The three disjoint
blocks in $R$ give
\[
 s_2(R)=1+s_2(S_c+T_c)+\lambda-s_2(S_c)
       =\lambda+s_2(T_c)+1-\tau_2(S_c,T_c),
\]
where $s_2$ is binary digit sum and
$\tau_2(S,T)=s_2(S)+s_2(T)-s_2(S+T)\ge0$ counts binary carries.
The strict inequality $S_c+T_c<L$ is essential.
The $31$ allowed low positions and the $22$ target positions give
\[
 s_2(T_c)=(W+23)\nu-31k-22,\qquad
 \lambda=1+k+E_*\nu.
\]
Therefore
\begin{align*}
 \lambda+s_2(T_c)+1
 &=5947\kappa+11864k-20\\
 &=2\bigl(15\kappa/2-10+2966(\kappa+2k)\bigr)
 =\log_2N^2.
\end{align*}
Legendre's formula yields
$v_2\binom{2R}{R}=s_2(R)$, hence
\[
 N^2\mid\binom{2R}{R}
 \quad\Longleftrightarrow\quad S_c\mathbin{\&}T_c=0.
\]
The small factor $2b$ in $L$ calibrates the threshold to a polynomial square
without an extra variable, an extra target mask, or a cubic radix.

\subsection{The interval guard and the Jones interface}
If $\Delta>0$ and $b\ge2$, then $K\ne0$ and
$0<g<c$. The identity $4g(c-g)=c^2-(c-2g)^2\le c^2$ implies
$4\delta^2<K^2$. Along with $s\ge1$, $\phi\ge0$, and $w\ne0$, this gives
$C,Y>0$ and $K>0$. The usual square-gap argument forces $h\ge1$:
if $h\le0$, then $1\le K\le R+1<P$; $K=1$ would force $P-1\mid R$,
which is impossible, while for $2\le K<P$,
\[
 (PK-1)^2<(P^2-1)K^2+1<(PK)^2.
\]
Thus $a_0$ cannot be a square in that branch.

For precision, instantiate the free ceiling in
\texttt{Jones1982.Nine.ratio\_sound} by
$G_{\mathrm{aux}}=8gK^2+1$, rather than assuming its old margin condition.
The positive integer $K^2-4\delta^2$ then gives
\[
 (K^2-4\delta^2)G_{\mathrm{aux}}-8gK^2>0.
\]
This ceiling is a proof-side abbreviation; the ratio registers do not
otherwise depend on it.

Replace $i$ by $|i|$, which changes no condition. To handle signed $j$,
$u\mid t$ and $F\equiv1\pmod J$ imply $F\mid H-C$. Also
$F\equiv1\pmod D$, $G\equiv1\pmod D$, and $G\equiv A\pmod F$ show
that the positive integers $D,F,I$ are pairwise coprime.
Hence $a_1=DFI$ square makes each a square. Put $d=G^2-1>0$ and
$x=\sqrt I$. The norm-one unit $x+H\sqrt d$ is positive even if $H<0$.
A positive power of $G+\sqrt d$ is the identity modulo $CF$; multiplying by
sufficiently large such powers gives a Pell solution $H'>B_1$ with
$H'\equiv H\pmod{CF}$. Then $j'=(H'-B_1)/C>0$, all square conditions
remain true, and $u\mid t$ remains true. The literal positive Jones interface
therefore applies and gives the power-of-two base and central-binomial
condition required above.

For completeness, take any common source zero, choose
$b=2^{1+2r}$ large enough, and set $\epsilon=2^r$. The direct source code
has $32g<G_c$, so the literal ratio completeness theorem supplies its
canonical witnesses. In \texttt{Ratio.lean}, \texttt{exists\_B25core} sets
\[
 \omega=2^{2R+1-k},\qquad U=(N^2/b)2^{2R+1}.
\]
Since $k=1+2r$ is odd and $k<2R+1$, this $\omega$ is a positive square.
No claim that the canonical $s$ is a square is needed.

The same canonical construction has
$\rho_0=(U+1)^{2R}/U^R=Y+f$ with
$0\le f<2^{2R}/U=b/(2N^2)$ and $Y\ge1$.
The source's \texttt{ratio\_core} gives
\[
 |C/K-\rho_0|\le\frac{2(Y+f)}{Y(U+1)}<\frac4U,
 \qquad |C/K-Y|<\frac b{N^2}.
\]
The last inequality uses $R\ge8$. Since $c,b\le N$ and
$g(c-g)\ge1$, it follows that
\[
 \frac{c^2(C-KY)^2}{K^2}<\frac{c^2b^2}{N^4}\le1\le g(c-g),
\]
so the very same canonical witnesses satisfy $\Delta>0$.
This is a cofinal square-base proof, not an unsupported restriction of an
existing slack variable to squares.

\subsection{Seeded square recovery, with all zero branches excluded}
For every integer assignment, $J=4A-5\ne0$ and $F\ne0$.
Indeed $F=0$ forces $A=0$ and $E=\pm1$, whereas $E$ is then divisible
by $J=-5$. Thus $u=JF\ne0$ unconditionally.
The polynomial
\[
 \Pi(x)=(x^3-x-1)(x^2+23)
\]
has no rational root and has a root modulo every positive integer.
For an odd prime other than $23$, either its cubic has a simple root or
that cubic is irreducible, whose Frobenius three-cycle makes its
nonzero discriminant $-23$ a square; the quadratic then has a simple root.
At $23$ the cubic has the simple root $3$. At $2$, $-23\equiv1\pmod8$
has compatible square roots. Hensel lifting and CRT finish the assertion.
Consequently $\Theta=u^5\Pi(t/u+r)\ne0$ for every integer assignment.
Moreover
\[
 \exists r\quad Zu^5\mid\Theta
 \quad\Longleftrightarrow\quad Z\ne0\ \text{and}\ u\mid t.
\]
For the reverse implication, choose a root of $\Pi$ modulo $|Z|$ after
$u\mid t$. For the forward implication, if a prime has
$v_p(u)=\alpha>\beta=v_p(t)$, the two factors of $\Theta$ have valuations
$3\beta$ and $2\beta$, contradicting $u^5\mid\Theta$.
If $Z=0$, divisibility would demand $\Theta=0$, already excluded.

At a zero of \eqref{eq:jones-seeded}, $\Sigma=0$ is impossible because
then $\mathcal F_{13}=\Theta^4\ne0$. Therefore
$\epsilon,w,i,z_1,z_2,z_3,a_0$ are all nonzero.
Every $d_j\equiv2\pmod3$ is nonzero, and every
$B_j^*\equiv1\pmod4$ is nonzero. Furthermore $B_j^*\ne1,-1$, so
$|B_j^*|\ge3$ and $|C_*|>1$ without any positivity hypothesis.

Let $q=X/\Sigma$. Equation \eqref{eq:jones-seeded} gives the exact identity
\[
 (q^2+C_*-C_*^2a_1)^2=4C_*q^2.
\]
If $q=0$, it would force $C_*a_1=1$, impossible since $|C_*|>1$.
Thus
\[
 r_0=\frac{q^2+C_*-C_*^2a_1}{2q},\qquad
 r_1=\frac{q-r_0}{C_*}
\]
are rational and satisfy $r_0^2=C_*$, $r_1^2=a_1$.
Rational square roots of integers are integers. Hence $C_*,a_1$ are
integer squares, $q\in\mathbb Z$, and $\Sigma\mid\Theta$.
This proof covers negative radicands and all possible cancellations without
an implicit positivity or distinct-square-class assumption.

Now $\gcd(B_j^*,C_{j-1})=1$, so the nonzero factors
$a_0,B_1^*,B_2^*,B_3^*$ are pairwise coprime. Since their product $C_*$
is a square, every absolute factor is a square. Each $B_j^*\equiv1\pmod4$
is therefore positive, because a negative square cannot be $1$ modulo $4$.
The sign of the product forces $a_0>0$, and $a_0$ is a square too.
All prefixes are now positive. A negative $d_j$ would make
$B_j^*\le-3$, so all $d_j>0$, yielding $s\ge1$, $\phi\ge0$, and
$\Delta>0$. This argument does not assert that $a_0\equiv1\pmod4$;
that false shortcut would invalidate the seeded construction.

Conversely, suppose $a_0$ is a positive square and these three sign
conditions hold. Each $d_j$ is positive and nonsquare since it is $2$
modulo $3$. Inductively $C_{j-1}$ is a positive square, so
$d_jC_{j-1}$ is positive nonsquare. A nontrivial Pell solution
$x^2-d_jC_{j-1}y^2=1$ can be chosen with $y$ even; set $z_j=y/2\ne0$.
Then $B_j^*=x^2$ and $C_j$ is a positive square. Given also $a_1$ square
and $u\mid t$, the intersective lemma chooses $r$ with
$\Sigma\mid\Theta$, and an integer choice of $v$ makes
$X/\Sigma=\sqrt{C_*}+C_*\sqrt{a_1}$.
Thus \eqref{eq:jones-seeded} vanishes. Combining both directions proves
the asserted equivalence to the original source and hence to a finite ZFC
refutation.

\subsection{Exact degree, circuit accounting, and provenance}
For the seeded equation the relevant exact register degrees are
\[
\begin{array}{c|rrrrrr}
 &R&N&a_0&a_1&C_*&\Sigma\\\hline
 \deg&45932&11864&652210&4363426&5592732&6692101
\end{array}
\]
and $\deg X=6692102$, $\deg A_0=18976934$,
$\deg A_1=28933092$. The last dominates $X^2$ and $A_0$, so the
output's leading term is $A_1^2$. The symbolic diagonal certificate has
leading coefficient
\[
 356960\,a^{14465376}\eta^{1888}\pmod{1000003},
\]
which is nonzero because $a$ and $\eta$ are powers of two.
The circuit upper bound and this surviving diagonal coefficient prove
$\deg\mathcal F_{13}=2(28933092)=57866184$ exactly, not merely at most.

The metric counts every live binary addition, subtraction and multiplication,
including scalar multiplications and every gate in fixed exponentiation.
Fixed coefficient leaves and reuse are free; coefficient bit size is a
separate measure.

The unseeded three-root baseline, also included, uses the same source,
starts $C_0=1$, omits $a_0$ from $Z$, and combines the square conditions
$C_*,a_0,a_1$ with the established separating weights
$1,C_*,C_*(P^2+K^2+2)$. It has $770$ operations and degree $71382096$.
The seeded improvement saves eleven operations and reduces this degree by
$18.93\%$ without changing the thirteen unknowns.
For completeness, if $W_2=C_*(P^2+K^2+2)$, its radicands are
$A_0=\Sigma^2C_*$, $A_1=\Sigma^2C_*^2a_0$ and
$A_2=\Sigma^2W_2^2a_1$. It uses
\[
 V=X^2+A_2-A_0-A_1,\quad
 W'=V^2+4X^2A_2-4A_0A_1,\quad
 (W')^2-16V^2X^2A_2.
\]
This is the product over all eight signs of
$X-\Sigma(\pm\sqrt{C_*}\pm C_*\sqrt{a_0}\pm W_2\sqrt{a_1})$.
The inequalities $|C_*|>\sqrt{|C_*|}$ and
$\sqrt{|a_0|}\le P^2+K^2+1$ make each weighted nonzero radical larger
in complex absolute value than the sum of all preceding radicals.
Thus effective sign patterns are distinct, also when a radicand is zero.
A rational signed sum is fixed by every multiquadratic automorphism;
uniqueness forces each nonzero root rational and hence integral.
This recovers all three square conditions and the merger divisibility.
At $\Sigma=0$ its output is $\Theta^8\ne0$. The unseeded Pell stack and
all remaining source and ratio steps are the same as proved above.

Run \texttt{python build\_jones.py} and
\texttt{python check\_jones.py}. The checker authenticates the fixed payloads
and pinned sources, independently replays the literal registers over three
prime fields (including signed and zero inputs), verifies the formal norm
identities, checks $136000$ signed seeded stacks, and exercises nonzero
merger and sparse-borrow branches. These finite tests corroborate the
universal proofs, not replace them.
All imported source paths, Git blob IDs, and SHA256 values are in
\texttt{SOURCE\_PINS.json}, at repository commit
\texttt{124c07c0b77a9e20e84676514d8b1ff78536ea3f}.
The inherited compiler and Jones ratio theorems are explicitly identified;
no newly compiled kernel certificate for this concrete equation is claimed.
